\documentclass[11pt,a4paper]{article}
\usepackage[utf8]{inputenc}
\usepackage[T1]{fontenc}
\usepackage[french]{babel}
\usepackage{amsmath}
\usepackage{amssymb}
\usepackage{geometry}
\usepackage{enumitem}
\usepackage{xcolor}
\usepackage{titlesec}
\usepackage{fancyhdr}
\usepackage{booktabs}
\usepackage{array}
\usepackage{longtable}
\usepackage{listings}
\usepackage{tcolorbox}
\usepackage{multirow}

% Configuration de la page
\geometry{margin=2.5cm}
\pagestyle{fancy}
\fancyhf{}
\fancyhead[L]{\leftmark}
\fancyhead[R]{NoSQL - CAP, ACID, CRUD}
\fancyfoot[C]{\thepage}

% Configuration des titres
\titleformat{\section}
{\Large\bfseries\color{blue!70!black}}
{}
{0em}
{}[\titlerule]

\titleformat{\subsection}
{\large\bfseries\color{blue!50!black}}
{}
{0em}
{}

\titleformat{\subsubsection}
{\normalsize\bfseries\color{blue!30!black}}
{}
{0em}
{}

% Boîtes colorées
\newtcolorbox{definitionbox}{
    colback=blue!5!white,
    colframe=blue!75!black,
    title=\textbf{Définition},
    fonttitle=\bfseries
}

\newtcolorbox{exemplebox}{
    colback=green!5!white,
    colframe=green!75!black,
    title=\textbf{Exemple},
    fonttitle=\bfseries
}

% Configuration pour le code
\lstset{
    basicstyle=\ttfamily\small,
    keywordstyle=\color{blue}\bfseries,
    commentstyle=\color{gray},
    stringstyle=\color{red},
    numbers=left,
    numberstyle=\tiny\color{gray},
    stepnumber=1,
    numbersep=5pt,
    backgroundcolor=\color{gray!10},
    frame=single,
    breaklines=true,
    showstringspaces=false,
    tabsize=2
}

% Métadonnées
\title{Résumé Complet NoSQL\\
\large CAP, ACID, CRUD pour Cassandra, MongoDB, Neo4j et Redis}
\author{AmineGR03}
\date{\today}

\begin{document}

\maketitle

\tableofcontents
\newpage

\section{Introduction}

Ce document présente les concepts fondamentaux des bases de données NoSQL et les opérations CRUD pour les principales technologies : Cassandra, MongoDB, Neo4j et Redis.

\newpage

\section{Théorème CAP}

\subsection{Définition}

Le \textbf{théorème CAP} (Consistency, Availability, Partition tolerance) énonce qu'un système distribué ne peut garantir simultanément que deux des trois propriétés suivantes :

\begin{definitionbox}
\textbf{Consistency (Cohérence)} : Tous les nœuds voient les mêmes données au même moment. Après une mise à jour, tous les nœuds retournent immédiatement la même valeur.

\textbf{Availability (Disponibilité)} : Le système reste opérationnel et répond aux requêtes même en cas de défaillance de certains nœuds.

\textbf{Partition tolerance (Tolérance aux pannes)} : Le système continue de fonctionner même si la communication entre certains nœuds est interrompue (partition réseau).
\end{definitionbox}

\subsection{Choix CAP par Base de Données}

\begin{table}[h]
\centering
\begin{tabular}{|l|c|c|c|}
\hline
\textbf{Base de données} & \textbf{C} & \textbf{A} & \textbf{P} \\
\hline
Cassandra & \textcolor{red}{$\times$} & \checkmark & \checkmark \\
\hline
MongoDB & \checkmark & \textcolor{red}{$\times$} & \checkmark \\
\hline
Neo4j & \checkmark & \textcolor{red}{$\times$} & \checkmark \\
\hline
Redis & \checkmark & \checkmark & \textcolor{red}{$\times$} \\
\hline
\end{tabular}
\caption{Choix CAP des bases de données NoSQL}
\end{table}

\subsection{Explications Détaillées}

\subsubsection{Cassandra : AP (Availability + Partition tolerance)}

\begin{itemize}
    \item \textbf{Architecture} : Système distribué avec réplication
    \item \textbf{Choix} : Privilégie la disponibilité et la tolérance aux pannes
    \item \textbf{Cohérence} : Cohérence éventuelle (Eventual Consistency)
    \item \textbf{Avantages} : Haute disponibilité, tolérance aux pannes réseau
    \item \textbf{Inconvénients} : Les lectures peuvent retourner des données obsolètes temporairement
    \item \textbf{Niveau de cohérence} : Configurable (ONE, QUORUM, ALL)
\end{itemize}

\subsubsection{MongoDB : CP (Consistency + Partition tolerance)}

\begin{itemize}
    \item \textbf{Architecture} : Réplication avec primary/secondary
    \item \textbf{Choix} : Privilégie la cohérence et la tolérance aux pannes
    \item \textbf{Disponibilité} : En cas de partition, le primary peut être indisponible
    \item \textbf{Avantages} : Données toujours cohérentes, tolérance aux pannes
    \item \textbf{Inconvénients} : Peut être indisponible lors d'une partition réseau
    \item \textbf{Transactions} : Supporte les transactions multi-documents (depuis v4.0)
\end{itemize}

\subsubsection{Neo4j : CP (Consistency + Partition tolerance)}

\begin{itemize}
    \item \textbf{Architecture} : Cluster avec leader/follower
    \item \textbf{Choix} : Privilégie la cohérence et la tolérance aux pannes
    \item \textbf{Disponibilité} : En cas de partition, le leader peut être indisponible
    \item \textbf{Avantages} : Données toujours cohérentes, relations garanties
    \item \textbf{Inconvénients} : Peut être indisponible lors d'une partition
    \item \textbf{Transactions} : Supporte les transactions ACID
\end{itemize}

\subsubsection{Redis : CA (Consistency + Availability)}

\begin{itemize}
    \item \textbf{Architecture} : Principalement en mémoire, peut être distribué
    \item \textbf{Choix} : Privilégie la cohérence et la disponibilité
    \item \textbf{Partition} : Redis Cluster supporte la partition tolerance
    \item \textbf{Avantages} : Données cohérentes, très rapide, haute disponibilité
    \item \textbf{Inconvénients} : Limité par la mémoire, partition tolerance optionnelle
    \item \textbf{Mode cluster} : Redis Cluster offre partition tolerance avec cohérence éventuelle
\end{itemize}

\newpage

\section{Propriétés ACID}

\subsection{Définition}

Les propriétés \textbf{ACID} garantissent la fiabilité des transactions dans une base de données :

\begin{definitionbox}
\textbf{Atomicity (Atomicité)} : Une transaction est "tout ou rien". Soit toutes les opérations réussissent, soit aucune n'est appliquée.

\textbf{Consistency (Cohérence)} : Une transaction amène la base de données d'un état cohérent à un autre état cohérent. Les contraintes d'intégrité sont respectées.

\textbf{Isolation (Isolation)} : Les transactions concurrentes ne se voient pas mutuellement. Chaque transaction s'exécute comme si elle était la seule.

\textbf{Durability (Durabilité)} : Une fois une transaction validée (committed), ses modifications persistent même en cas de panne système.
\end{definitionbox}

\subsection{Support ACID par Base de Données}

\begin{table}[h]
\centering
\begin{tabular}{|l|c|c|c|c|}
\hline
\textbf{Base de données} & \textbf{A} & \textbf{C} & \textbf{I} & \textbf{D} \\
\hline
Cassandra & \checkmark & \textcolor{orange}{Partiel} & \textcolor{orange}{Partiel} & \checkmark \\
\hline
MongoDB & \checkmark & \checkmark & \checkmark & \checkmark \\
\hline
Neo4j & \checkmark & \checkmark & \checkmark & \checkmark \\
\hline
Redis & \checkmark & \checkmark & \textcolor{orange}{Partiel} & \textcolor{orange}{Partiel} \\
\hline
\end{tabular}
\caption{Support ACID des bases de données NoSQL}
\end{table}

\subsection{Explications Détaillées}

\subsubsection{Cassandra et ACID}

\begin{itemize}
    \item \textbf{Atomicity} : \checkmark Supportée au niveau de la partition (BATCH)
    \item \textbf{Consistency} : \textcolor{orange}{Partielle} - Cohérence éventuelle par défaut, configurable
    \item \textbf{Isolation} : \textcolor{orange}{Partielle} - Isolation au niveau de la partition
    \item \textbf{Durability} : \checkmark Écritures persistées sur disque
    \item \textbf{Transactions} : Pas de transactions multi-partitions traditionnelles
    \item \textbf{LWT (Lightweight Transactions)} : Supporte les opérations conditionnelles atomiques
\end{itemize}

\subsubsection{MongoDB et ACID}

\begin{itemize}
    \item \textbf{Atomicity} : \checkmark Supportée au niveau du document (single-document) et multi-documents (depuis v4.0)
    \item \textbf{Consistency} : \checkmark Garantie par les transactions et les contraintes
    \item \textbf{Isolation} : \checkmark Niveaux d'isolation configurables (read uncommitted, read committed, snapshot)
    \item \textbf{Durability} : \checkmark Journaling (WAL) et réplication
    \item \textbf{Transactions} : Transactions multi-documents depuis MongoDB 4.0
    \item \textbf{Replica Set} : Garantit la cohérence entre réplicas
\end{itemize}

\subsubsection{Neo4j et ACID}

\begin{itemize}
    \item \textbf{Atomicity} : \checkmark Toutes les opérations dans une transaction sont atomiques
    \item \textbf{Consistency} : \checkmark Contraintes et règles de graphe garanties
    \item \textbf{Isolation} : \checkmark Isolation complète des transactions
    \item \textbf{Durability} : \checkmark Transactions journalisées et persistées
    \item \textbf{Transactions} : Support complet des transactions ACID
    \item \textbf{Graph Integrity} : Garantit l'intégrité des relations et nœuds
\end{itemize}

\subsubsection{Redis et ACID}

\begin{itemize}
    \item \textbf{Atomicity} : \checkmark Toutes les commandes sont atomiques, MULTI/EXEC pour transactions
    \item \textbf{Consistency} : \checkmark Garantie par l'exécution atomique des commandes
    \item \textbf{Isolation} : \textcolor{orange}{Partielle} - Pas d'isolation entre transactions (pas de verrous)
    \item \textbf{Durability} : \textcolor{orange}{Partielle} - Optionnelle (AOF, RDB), par défaut en mémoire
    \item \textbf{Transactions} : MULTI/EXEC pour exécution atomique (pas de rollback)
    \item \textbf{Persistence} : AOF (Append Only File) ou RDB (snapshots) pour durabilité
\end{itemize}

\newpage

\section{Opérations CRUD}

\subsection{Définition}

Les opérations \textbf{CRUD} sont les quatre opérations de base pour manipuler des données :

\begin{definitionbox}
\textbf{Create (Créer)} : Insérer de nouvelles données dans la base.

\textbf{Read (Lire)} : Récupérer des données existantes.

\textbf{Update (Mettre à jour)} : Modifier des données existantes.

\textbf{Delete (Supprimer)} : Supprimer des données de la base.
\end{definitionbox}

\newpage

\section{Cassandra - CRUD}

\subsection{Create (INSERT)}

\begin{lstlisting}[language=SQL]
-- Insertion simple
INSERT INTO utilisateurs (id, nom, email, age)
VALUES (uuid(), 'Dupont', 'dupont@email.com', 30);

-- Insertion avec TTL (Time To Live)
INSERT INTO utilisateurs (id, nom, email)
VALUES (uuid(), 'Martin', 'martin@email.com')
USING TTL 3600;

-- Insertion conditionnelle
INSERT INTO utilisateurs (id, nom, email)
VALUES (a1b2c3d4-e5f6-7890-abcd-ef1234567890, 'Dupont', 'dupont@email.com')
IF NOT EXISTS;
\end{lstlisting}

\subsection{Read (SELECT)}

\begin{lstlisting}[language=SQL]
-- Sélection de tous les champs
SELECT * FROM utilisateurs;

-- Sélection de champs spécifiques
SELECT id, nom, email FROM utilisateurs;

-- Sélection avec WHERE
SELECT * FROM utilisateurs WHERE id = a1b2c3d4-e5f6-7890-abcd-ef1234567890;

-- Sélection avec LIMIT
SELECT * FROM utilisateurs LIMIT 10;

-- Sélection avec COUNT
SELECT COUNT(*) FROM utilisateurs;

-- Sélection avec ALLOW FILTERING (attention: coûteux)
SELECT * FROM utilisateurs WHERE ville = 'Paris' AND age > 25
ALLOW FILTERING;
\end{lstlisting}

\subsection{Update (UPDATE)}

\begin{lstlisting}[language=SQL]
-- Mise à jour simple
UPDATE utilisateurs
SET age = 31, ville = 'Lyon'
WHERE id = a1b2c3d4-e5f6-7890-abcd-ef1234567890;

-- Mise à jour avec TTL
UPDATE utilisateurs
USING TTL 7200
SET email = 'nouveau@email.com'
WHERE id = a1b2c3d4-e5f6-7890-abcd-ef1234567890;

-- Mise à jour conditionnelle
UPDATE utilisateurs
SET age = 32
WHERE id = a1b2c3d4-e5f6-7890-abcd-ef1234567890
IF age = 31;

-- Mise à jour avec IF EXISTS
UPDATE utilisateurs
SET age = 33
WHERE id = a1b2c3d4-e5f6-7890-abcd-ef1234567890
IF EXISTS;
\end{lstlisting}

\subsection{Delete (DELETE)}

\begin{lstlisting}[language=SQL]
-- Suppression d'un document
DELETE FROM utilisateurs
WHERE id = a1b2c3d4-e5f6-7890-abcd-ef1234567890;

-- Suppression de champs spécifiques
DELETE email, age FROM utilisateurs
WHERE id = a1b2c3d4-e5f6-7890-abcd-ef1234567890;

-- Suppression conditionnelle
DELETE FROM utilisateurs
WHERE id = a1b2c3d4-e5f6-7890-abcd-ef1234567890
IF age = 30;

-- TRUNCATE (vider une table)
TRUNCATE utilisateurs;
\end{lstlisting}

\subsection{Tableau Récapitulatif CRUD Cassandra}

\begin{table}[h]
\centering
\begin{tabular}{|c|c|}
\hline
\textbf{Opération} & \textbf{Commande} \\
\hline
Create & \texttt{INSERT INTO ... VALUES ...} \\
\hline
Read & \texttt{SELECT ... FROM ... WHERE ...} \\
\hline
Update & \texttt{UPDATE ... SET ... WHERE ...} \\
\hline
Delete & \texttt{DELETE FROM ... WHERE ...} \\
\hline
\end{tabular}
\caption{Opérations CRUD Cassandra}
\end{table}

\newpage

\section{MongoDB - CRUD}

\subsection{Create (Insert)}

\begin{lstlisting}[language=JavaScript]
// Insertion d'un document
db.collection.insertOne({
    nom: "Dupont",
    age: 30,
    ville: "Paris"
})

// Insertion de plusieurs documents
db.collection.insertMany([
    {nom: "Martin", age: 25, ville: "Lyon"},
    {nom: "Bernard", age: 35, ville: "Marseille"}
])

// Insertion avec insert() (déprécié mais fonctionne)
db.collection.insert({
    nom: "Durand",
    age: 28,
    ville: "Toulouse"
})
\end{lstlisting}

\subsection{Read (Find)}

\begin{lstlisting}[language=JavaScript]
// Lecture de tous les documents
db.collection.find()

// Lecture avec formatage
db.collection.find().pretty()

// Lecture d'un seul document
db.collection.findOne({age: 30})

// Lecture avec filtre
db.collection.find({ville: "Paris"})

// Projection (sélection de champs)
db.collection.find({}, {nom: 1, age: 1, _id: 0})

// Lecture avec opérateurs
db.collection.find({age: {$gt: 25}})
db.collection.find({ville: {$in: ["Paris", "Lyon"]}})
\end{lstlisting}

\subsection{Update (Update)}

\begin{lstlisting}[language=JavaScript]
// Mise à jour d'un document
db.collection.updateOne(
    {nom: "Dupont"},
    {$set: {age: 31}}
)

// Mise à jour de plusieurs documents
db.collection.updateMany(
    {ville: "Paris"},
    {$set: {pays: "France"}}
)

// Opérateurs de mise à jour
{$set: {age: 30}}        // Définir une valeur
{$inc: {age: 1}}         // Incrémenter
{$unset: {ville: ""}}   // Supprimer un champ
{$push: {hobbies: "lecture"}}  // Ajouter à un tableau
{$pull: {hobbies: "lecture"}}  // Retirer d'un tableau

// Replace (remplacement complet)
db.collection.replaceOne(
    {nom: "Dupont"},
    {nom: "Dupont", age: 32, ville: "Lyon"}
)
\end{lstlisting}

\subsection{Delete (Delete)}

\begin{lstlisting}[language=JavaScript]
// Suppression d'un document
db.collection.deleteOne({nom: "Dupont"})

// Suppression de plusieurs documents
db.collection.deleteMany({ville: "Paris"})

// Suppression de tous les documents
db.collection.deleteMany({})
\end{lstlisting}

\subsection{Tableau Récapitulatif CRUD MongoDB}

\begin{table}[h]
\centering
\begin{tabular}{|c|c|}
\hline
\textbf{Opération} & \textbf{Commande} \\
\hline
Create & \texttt{insertOne()}, \texttt{insertMany()} \\
\hline
Read & \texttt{find()}, \texttt{findOne()} \\
\hline
Update & \texttt{updateOne()}, \texttt{updateMany()} \\
\hline
Delete & \texttt{deleteOne()}, \texttt{deleteMany()} \\
\hline
\end{tabular}
\caption{Opérations CRUD MongoDB}
\end{table}

\newpage

\section{Neo4j - CRUD}

\subsection{Create (CREATE, MERGE)}

\begin{lstlisting}[language=SQL]
-- Créer un nœud
CREATE (p:Person {nom: 'Dupont', age: 30, ville: 'Paris'})
RETURN p;

-- Créer plusieurs nœuds
CREATE 
    (p1:Person {nom: 'Dupont', age: 30}),
    (p2:Person {nom: 'Martin', age: 25}),
    (p3:Person {nom: 'Bernard', age: 35})
RETURN p1, p2, p3;

-- Créer une relation
CREATE (p1:Person {nom: 'Dupont'})-[:CONNAIT]->(p2:Person {nom: 'Martin'})
RETURN p1, p2;

-- MERGE (créer si n'existe pas, trouver si existe)
MERGE (p:Person {nom: 'Dupont'})
ON CREATE SET p.age = 30, p.ville = 'Paris'
ON MATCH SET p.age = 31
RETURN p;
\end{lstlisting}

\subsection{Read (MATCH)}

\begin{lstlisting}[language=SQL]
-- MATCH simple
MATCH (p:Person)
RETURN p;

-- MATCH avec filtre
MATCH (p:Person {nom: 'Dupont'})
RETURN p;

-- MATCH avec WHERE
MATCH (p:Person)
WHERE p.age > 25
RETURN p;

-- MATCH avec relations
MATCH (p1:Person)-[:CONNAIT]->(p2:Person)
RETURN p1, p2;

-- MATCH avec profondeur variable
MATCH (p:Person)-[:CONNAIT*1..3]->(autre:Person)
RETURN p, autre;

-- OPTIONAL MATCH (jointure externe)
MATCH (p:Person)
OPTIONAL MATCH (p)-[:CONNAIT]->(ami:Person)
RETURN p, ami;
\end{lstlisting}

\subsection{Update (SET)}

\begin{lstlisting}[language=SQL]
-- SET simple
MATCH (p:Person {nom: 'Dupont'})
SET p.age = 31
RETURN p;

-- SET multiple
MATCH (p:Person {nom: 'Dupont'})
SET p.age = 31, p.ville = 'Lyon'
RETURN p;

-- Ajouter un label
MATCH (p:Person {nom: 'Dupont'})
SET p:Client
RETURN p;

-- Mettre à jour une relation
MATCH (p1:Person {nom: 'Dupont'})-[r:CONNAIT]->(p2:Person {nom: 'Martin'})
SET r.depuis = 2021
RETURN r;
\end{lstlisting}

\subsection{Delete (DELETE, DETACH DELETE, REMOVE)}

\begin{lstlisting}[language=SQL]
-- DELETE nœud (nécessite de supprimer les relations d'abord)
MATCH (p:Person {nom: 'Dupont'})
DELETE p;

-- DETACH DELETE (supprime nœud et toutes ses relations)
MATCH (p:Person {nom: 'Dupont'})
DETACH DELETE p;

-- DELETE relation
MATCH (p1:Person {nom: 'Dupont'})-[r:CONNAIT]->(p2:Person {nom: 'Martin'})
DELETE r;

-- REMOVE (supprimer propriété ou label)
MATCH (p:Person {nom: 'Dupont'})
REMOVE p.ville
RETURN p;

MATCH (p:Person {nom: 'Dupont'})
REMOVE p:Client
RETURN p;
\end{lstlisting}

\subsection{Tableau Récapitulatif CRUD Neo4j}

\begin{table}[h]
\centering
\begin{tabular}{|c|c|}
\hline
\textbf{Opération} & \textbf{Commande} \\
\hline
Create & \texttt{CREATE}, \texttt{MERGE} \\
\hline
Read & \texttt{MATCH ... RETURN} \\
\hline
Update & \texttt{MATCH ... SET} \\
\hline
Delete & \texttt{MATCH ... DELETE}, \texttt{DETACH DELETE} \\
\hline
Remove & \texttt{MATCH ... REMOVE} \\
\hline
\end{tabular}
\caption{Opérations CRUD Neo4j}
\end{table}

\newpage

\section{Redis - CRUD}

\subsection{Create (SET, HSET, LPUSH, RPUSH, SADD, ZADD)}

\begin{lstlisting}
-- String
SET cle "valeur"
SET nom "Dupont"
SET age 30

-- Hash
HSET utilisateur:1 nom "Dupont"
HSET utilisateur:1 age 30
HSET utilisateur:1 ville "Paris"

-- List
LPUSH liste "element1"
RPUSH liste "element2"

-- Set
SADD tags "important"
SADD tags "urgent" "prioritaire"

-- Sorted Set
ZADD scores 100 "joueur1"
ZADD scores 200 "joueur2" 150 "joueur3"
\end{lstlisting}

\subsection{Read (GET, HGET, LRANGE, SMEMBERS, ZRANGE)}

\begin{lstlisting}
-- String
GET cle
GET nom

-- Hash
HGET utilisateur:1 nom
HGETALL utilisateur:1

-- List
LRANGE liste 0 -1
LINDEX liste 0

-- Set
SMEMBERS tags
SISMEMBER tags "important"

-- Sorted Set
ZRANGE scores 0 -1
ZRANGE scores 0 -1 WITHSCORES
ZREVRANGE scores 0 -1
ZSCORE scores "joueur1"
\end{lstlisting}

\subsection{Update (SET, HSET, LSET, ZADD)}

\begin{lstlisting}
-- String (remplace la valeur)
SET nom "Martin"

-- Hash
HSET utilisateur:1 age 31

-- List
LSET liste 0 "nouvelle_valeur"

-- Set (ajouter)
SADD tags "nouveau_tag"

-- Sorted Set (modifier score)
ZADD scores 250 "joueur1"
ZINCRBY scores 10 "joueur1"
\end{lstlisting}

\subsection{Delete (DEL, HDEL, LPOP, RPOP, SREM, ZREM)}

\begin{lstlisting}
-- String
DEL cle
DEL nom age

-- Hash
HDEL utilisateur:1 ville
HDEL utilisateur:1 nom age

-- List
LPOP liste
RPOP liste
LREM liste 2 "val"

-- Set
SREM tags "important"
SREM tags "urgent" "prioritaire"

-- Sorted Set
ZREM scores "joueur1"
ZREM scores "joueur2" "joueur3"
\end{lstlisting}

\subsection{Tableau Récapitulatif CRUD Redis}

\begin{table}[h]
\centering
\begin{tabular}{|c|c|c|c|c|}
\hline
\textbf{Type} & \textbf{Create} & \textbf{Read} & \textbf{Update} & \textbf{Delete} \\
\hline
String & \texttt{SET} & \texttt{GET} & \texttt{SET} & \texttt{DEL} \\
\hline
Hash & \texttt{HSET} & \texttt{HGET}, \texttt{HGETALL} & \texttt{HSET} & \texttt{HDEL} \\
\hline
List & \texttt{LPUSH/RPUSH} & \texttt{LRANGE} & \texttt{LSET} & \texttt{LPOP/RPOP} \\
\hline
Set & \texttt{SADD} & \texttt{SMEMBERS} & \texttt{SADD} & \texttt{SREM} \\
\hline
Sorted Set & \texttt{ZADD} & \texttt{ZRANGE} & \texttt{ZADD} & \texttt{ZREM} \\
\hline
\end{tabular}
\caption{Opérations CRUD Redis par type de données}
\end{table}

\newpage

\section{Comparaison CRUD entre Bases de Données}

\subsection{Tableau Comparatif}

\begin{table}[h]
\centering
\small
\begin{tabular}{|l|p{3cm}|p{3cm}|p{3cm}|p{3cm}|}
\hline
\textbf{Opération} & \textbf{Cassandra} & \textbf{MongoDB} & \textbf{Neo4j} & \textbf{Redis} \\
\hline
\multirow{2}{*}{Create} & \texttt{INSERT INTO} & \texttt{insertOne()} & \texttt{CREATE} & \texttt{SET/HSET/} \\
& \texttt{VALUES} & \texttt{insertMany()} & \texttt{MERGE} & \texttt{LPUSH/SADD/ZADD} \\
\hline
\multirow{2}{*}{Read} & \texttt{SELECT ...} & \texttt{find()} & \texttt{MATCH ...} & \texttt{GET/HGET/} \\
& \texttt{FROM ... WHERE} & \texttt{findOne()} & \texttt{RETURN} & \texttt{LRANGE/SMEMBERS} \\
\hline
\multirow{2}{*}{Update} & \texttt{UPDATE ...} & \texttt{updateOne()} & \texttt{MATCH ... SET} & \texttt{SET/HSET/} \\
& \texttt{SET ... WHERE} & \texttt{updateMany()} & & \texttt{LSET/ZADD} \\
\hline
\multirow{2}{*}{Delete} & \texttt{DELETE FROM} & \texttt{deleteOne()} & \texttt{MATCH ...} & \texttt{DEL/HDEL/} \\
& \texttt{WHERE} & \texttt{deleteMany()} & \texttt{DELETE} & \texttt{LPOP/SREM/ZREM} \\
\hline
\end{tabular}
\caption{Comparaison des opérations CRUD}
\end{table}

\subsection{Caractéristiques Spécifiques}

\subsubsection{Cassandra}

\begin{itemize}
    \item Syntaxe SQL-like (CQL)
    \item Requêtes limitées par la structure de la clé primaire
    \item Support des opérations conditionnelles (IF EXISTS, IF NOT EXISTS)
    \item TTL (Time To Live) pour expiration automatique
    \item BATCH pour opérations atomiques
\end{itemize}

\subsubsection{MongoDB}

\begin{itemize}
    \item Syntaxe JavaScript/JSON
    \item Documents flexibles (schéma dynamique)
    \item Opérateurs de mise à jour riches (\$set, \$inc, \$push, etc.)
    \item Projection pour sélectionner des champs
    \item Agrégation pipeline pour requêtes complexes
\end{itemize}

\subsubsection{Neo4j}

\begin{itemize}
    \item Syntaxe Cypher (déclarative)
    \item Focus sur les relations entre nœuds
    \item MERGE pour créer ou mettre à jour
    \item DETACH DELETE pour supprimer nœuds et relations
    \item Requêtes de graphe (chemins, profondeur variable)
\end{itemize}

\subsubsection{Redis}

\begin{itemize}
    \item Commandes simples et directes
    \item Structures de données variées (String, Hash, List, Set, Sorted Set)
    \item Opérations atomiques par commande
    \item Transactions avec MULTI/EXEC
    \item TTL pour expiration automatique
\end{itemize}

\newpage

\section{Résumé et Tableaux Récapitulatifs}

\subsection{Tableau CAP}

\begin{table}[h]
\centering
\begin{tabular}{|l|c|c|c|}
\hline
\textbf{Base de données} & \textbf{Choix CAP} & \textbf{Type} & \textbf{Caractéristique principale} \\
\hline
Cassandra & AP & Wide-column & Haute disponibilité, répartition \\
\hline
MongoDB & CP & Document & Cohérence forte, réplication \\
\hline
Neo4j & CP & Graph & Cohérence, relations garanties \\
\hline
Redis & CA* & Key-value & Performance, cohérence \\
\hline
\end{tabular}
\caption{Classification CAP des bases NoSQL}
\end{table}

*Redis Cluster offre partition tolerance

\subsection{Tableau ACID}

\begin{table}[h]
\centering
\begin{tabular}{|l|c|c|}
\hline
\textbf{Base de données} & \textbf{Support ACID} & \textbf{Notes} \\
\hline
Cassandra & Partiel & Atomicité par partition, cohérence éventuelle \\
\hline
MongoDB & Complet & Transactions multi-documents (v4.0+) \\
\hline
Neo4j & Complet & Transactions ACID complètes \\
\hline
Redis & Partiel & Atomicité par commande, durabilité optionnelle \\
\hline
\end{tabular}
\caption{Support ACID des bases NoSQL}
\end{table}

\subsection{Tableau CRUD - Vue d'Ensemble}

\begin{longtable}{|p{2cm}|p{3.5cm}|p{3.5cm}|p{3.5cm}|p{3.5cm}|}
\hline
\textbf{Opération} & \textbf{Cassandra} & \textbf{MongoDB} & \textbf{Neo4j} & \textbf{Redis} \\
\hline
\endhead
\multirow{2}{*}{Create} & \texttt{INSERT INTO} & \texttt{insertOne()} & \texttt{CREATE} & \texttt{SET} \\
& \texttt{... VALUES} & \texttt{insertMany()} & \texttt{MERGE} & \texttt{HSET/ZADD} \\
\hline
\multirow{2}{*}{Read} & \texttt{SELECT ...} & \texttt{find()} & \texttt{MATCH} & \texttt{GET} \\
& \texttt{FROM ... WHERE} & \texttt{findOne()} & \texttt{... RETURN} & \texttt{HGETALL} \\
\hline
\multirow{2}{*}{Update} & \texttt{UPDATE ...} & \texttt{updateOne()} & \texttt{MATCH ... SET} & \texttt{SET} \\
& \texttt{SET ... WHERE} & \texttt{updateMany()} & & \texttt{HSET} \\
\hline
\multirow{2}{*}{Delete} & \texttt{DELETE FROM} & \texttt{deleteOne()} & \texttt{MATCH ...} & \texttt{DEL} \\
& \texttt{... WHERE} & \texttt{deleteMany()} & \texttt{DELETE} & \texttt{HDEL} \\
\hline
\caption{Tableau comparatif complet CRUD}
\end{longtable}

\newpage

\section{Conseils et Bonnes Pratiques}

\subsection{Choix de Base de Données}

\begin{itemize}
    \item \textbf{Cassandra} : Données distribuées, haute disponibilité, écritures intensives
    \item \textbf{MongoDB} : Documents flexibles, requêtes complexes, évolutivité horizontale
    \item \textbf{Neo4j} : Relations complexes, graphes sociaux, recommandations
    \item \textbf{Redis} : Cache, sessions, queues, données temporaires, performance
\end{itemize}

\subsection{Considérations CAP}

\begin{itemize}
    \item \textbf{AP} : Accepte la cohérence éventuelle pour la disponibilité
    \item \textbf{CP} : Priorise la cohérence, peut sacrifier la disponibilité
    \item \textbf{CA} : Fonctionne bien en réseau local, limité en distribution
\end{itemize}

\subsection{Considérations ACID}

\begin{itemize}
    \item Les bases NoSQL offrent souvent un compromis entre ACID et performance
    \item MongoDB et Neo4j offrent un support ACID complet
    \item Cassandra et Redis offrent des garanties partielles mais optimisées
    \item Choisir selon les besoins de cohérence de l'application
\end{itemize}

\subsection{Opérations CRUD}

\begin{itemize}
    \item \textbf{Cassandra} : Attention aux requêtes (limitées par la clé primaire)
    \item \textbf{MongoDB} : Utiliser les index pour optimiser les requêtes
    \item \textbf{Neo4j} : Exploiter les relations pour des requêtes efficaces
    \item \textbf{Redis} : Choisir la bonne structure de données pour chaque cas d'usage
\end{itemize}

\vspace{2cm}

\begin{center}
\textbf{Fin du document}
\end{center}

\end{document}

